\documentclass[12pt]{article}
\newif\ifblind
\blindfalse
% \blindtrue

\input{preamble_paper.tex}
% Combined paper, formatted for the Journal of Economic Psychology (research article: max. 12,000 words including abstract,
% references, tables, figures and appendix; no footnotes on initial submission; p-values with three decimals).
% Compile from papers/:  latexmk -pdf -outdir=build wellbeing.tex 
% Submission strategy: see TODO.md, section 8.

% Is GDP a Poor Predictor of National Well-Being? Evidence from Two Surveys and a Wording Experiment
% Region Beats Income in Predicting National Well-Being, Unless You Ask Gallup
\title{Does Money Buy National Happiness?\\ It Depends on Who You Ask}
\ifblind
  \author{}
  \date{}
\else
    \renewcommand{\url}[1]{\href{#1}{Link}} 
    \author{Adrien Fabre\thanks{CNRS; CIRED. adrien.fabre@cnrs.fr}\,\thanks{Data and code: \href{https://github.com/bixiou/wellbeing_gdp_region}{github.com/bixiou/wellbeing\_gdp\_region}. %\todo{deposit data and code on a public repository (e.g.\ OSF) before acceptance; Gallup data cannot be redistributed}. 
    The survey was approved by the CIRED institutional review board (IRB-CIRED-2025-2) and pre-registered at \href{https://osf.io/7mzn4}{osf.io/7mzn4}. I thank Claudia Senik and participants in the 2024 STATEC Well-being Conference for helpful comments, Armon Rezai and the Vienna University of Economics and Business (WU Wien) for access to the Gallup World Poll data, and Julieta Toffoli for research assistance. This research received funding from Agence Nationale de la Recherche (ANR-24-CE03-7110). Declaration of interest: none.}}
    \date{\today}
    % \vspace{-1cm}
\fi 

\begin{document}
\maketitle

\begin{abstract}
\noindent Is GDP per capita a good predictor of national well-being? Using all waves of the World Values Survey and the European Values Study (WVS/EVS, 1981--2023, \nWVSCountries countries), I find that log GDP per capita explains on average \meanRtwoIncomePPP\% of the cross-country variance of eight well-being indicators, and that a country's broad world region is a better predictor, in and out of sample, in \shareRegionBetterWVS\% of \nSpecsWVS specifications. This advantage stems from Latin America and the former Eastern Bloc, and it vanishes in the Gallup World Poll, where income predicts life evaluations as well as region. To understand why the two datasets disagree, I randomize the question wording (Cantril ladder vs.\ life satisfaction) and scale (0--10 vs.\ 1--10) in an online survey of \nFabre respondents in ten high-income countries. The ladder shifts answers by \effectLadder points, with effects varying across countries, which explains why Gallup levels are lower than WVS levels. However, in the same ten countries, the two questions yield similar income gradients when asked to the same sample, whereas Gallup's gradient is three times steeper than the WVS's: the discrepancy stems from differences in samples, modes and context rather than wording.
\end{abstract}

\noindent\textbf{Keywords:} subjective well-being; life satisfaction; GDP; world regions; survey methodology; question wording.\\
\noindent\textbf{JEL codes:} I31; O15; O57; C83.

\ifblind
  \doublespacing
  \linenumbers
\else
  \tableofcontents
\fi

\clearpage

\section{Introduction}\label{sec:intro}

Which country is the happiest? The answer most often heard, based on the World Happiness Report, is a Scandinavian, high-income country. Behind this answer lies one of the most robust stylized facts of well-being research: across countries, average life evaluations increase with the logarithm of GDP per capita \citep{deaton_income_2008,stevenson_economic_2008,sacks_new_2012}. Using the Gallup World Poll, \citet{deaton_income_2008} finds a correlation above 0.8 between average life evaluation and log GDP per capita; using the World Values Survey (WVS), \citet{inglehart_genes_2000} reported an $R^2$ of about one half for an index of happiness and satisfaction. These correlations are frequently invoked to support the view that economic growth is a reliable route to higher well-being.

This paper revisits the predictive capacity of GDP per capita for national well-being, and asks two questions. First, how much of the cross-country variation in well-being does GDP per capita actually predict, and does it predict it better than a very simple alternative, the country's world region? Second, why do the two main sources of cross-national well-being data, Gallup and the WVS, give different answers to the first question? Previous work has noted that income is more strongly related to Gallup's ladder than to WVS measures \citep{deaton_income_2008}, and that national income is more related to life evaluations than to happiness or emotions \citep{inglehart_development_2008,kahneman_high_2010}. Yet whether the difference between datasets comes from the question asked or from the samples has, to my knowledge, not been tested experimentally.

For the first question, I use all seven waves of the WVS, complemented with the 2017 wave of the European Values Study (\nWVSCountryYears country-years), eight indicators of national well-being, and seven income measures, and I compare the variance explained by income and by region using the Shapley (LMG) decomposition of $R^2$ \citep{lindeman_introduction_1980,gromping_estimators_2007} and leave-one-country-out cross-validation. For the second question, I embedded a $2 \times 2$ experiment in an online survey of \nFabre respondents in ten high-income countries, randomizing the wording (Gallup's Cantril ladder vs.\ the WVS life-satisfaction question) and the scale (0--10 vs.\ 1--10). Combining the experiment with Gallup and WVS data for the same countries, I decompose the Gallup--WVS gap into a question effect and a residual reflecting differences in samples, modes, context and time.

The results are threefold. First, in the WVS, GDP per capita is a weak predictor of national well-being: it explains on average \meanRtwoIncomePPP\% of the variance of the eight indicators (\meanRtwoIncomeBest\% with the best income measure), and nothing of the share of very happy people. Region is a better predictor in \shareRegionBetterWVS\% of specifications and, out of sample, for every indicator. Second, this conclusion is fragile in two respects: it is driven by Latin America and the former Eastern Bloc, and it does not hold in Gallup, where income predicts life evaluations as well as or better than region: log GDP per capita explains \RtwoGallupMeanPPP\% of the cross-country variance of Gallup's mean ladder, against \RtwoGallupRegion\% for region. Third, the experiment shows that question wording has a sizeable effect on levels---the ladder shifts answers by \effectLadder points, more than the average gap between Gallup and the WVS---and that this effect varies across countries. However, wording has little effect on the relationship between income and well-being, which is what the two datasets disagree about. Across the ten countries, the slope of mean life evaluation on log GDP per capita is \slopeTenGallup in Gallup but \slopeTenWVS in the WVS, whereas in the experiment it is \slopeTenFabreLadder with the ladder and \slopeTenFabreSatisfaction with the satisfaction question. Likewise, within countries, richer respondents do not rate their lives relatively higher when asked the ladder question rather than the satisfaction question. The discrepancy between Gallup and the WVS is therefore attributable to differences between samples in the broad sense, not to the question asked.

The paper contributes to three literatures. It adds to the debate on the relationship between income and well-being across countries and over time \citep{easterlin_does_1974,easterlin_happiness_2010,stevenson_economic_2008,deaton_income_2008,jebb_happiness_2018} by showing that the strength of the cross-country relationship depends on the dataset, and by benchmarking income against regions. It relates to work on regional and cultural differences in well-being, in particular the high well-being of Latin America \citep{graham_paradox_2009,rojas_handbook_2016} and the low well-being of transition countries \citep{guriev_unhappiness_2009}, and on the comparability of well-being across cultures \citep{diener_culture_2000,exton_comparing_2015,senik_french_2014}. Finally, it contributes to the methodological literature on the measurement of subjective well-being \citep{kahneman_developments_2006,oecd_guidelines_2013,stone_subjective_2013,bond_sad_2019,kaiser_scientific_2022}, and specifically on the comparability of the Cantril ladder and life-satisfaction questions \citep{bjornskov_how_2010,nilsson_cantril_2024} and on context and mode effects \citep{deaton_understanding_2016,dolan_happy_2016}, by providing experimental estimates of wording and scale effects in ten countries.

Section~\ref{sec:data} presents the data. Section~\ref{sec:region} compares income and region as predictors of national well-being. Section~\ref{sec:discrepancy} investigates the discrepancy between Gallup and the WVS. Section~\ref{sec:discussion} discusses limitations, and Section~\ref{sec:conclusion} concludes.

\section{Data}\label{sec:data}
% Shared section: data on national well-being, income and regions (WVS, Gallup, WHR, WDI)
\subsection{World Values Survey}\label{subsec:wvs}

The World Values Survey (WVS) is the longest-running cross-national survey of values and subjective well-being. I use the WVS time-series file \citep{wvs_trend_2022}, which pools the seven waves conducted between 1981 and 2022, and complete it with the Joint EVS/WVS 2017--2022 dataset \citep{evs_wvs_joint_2024}, which adds the \nEVSSurveys national surveys of the 2017 wave of the European Values Study (EVS, conducted in 2017--2021) and \nNewWVSSurveys recent WVS surveys (India 2023, Uzbekistan 2022). The EVS asks the same well-being questions, and the combination of the two surveys is standard practice (the so-called Integrated Values Surveys). The EVS and the WVS are distinct surveys: when both surveyed the same country the same year (\nEVSSameYearWVS cases), I keep both surveys as separate observations.
For brevity, I refer to the combined data as the WVS: \nWVSRespondents respondents in \nWVSCountryYears country-years and \nWVSCountries countries. National samples are designed to be representative of the adult population, typically with 1,000 to 3,000 respondents, and most interviews are conducted face-to-face (though some recent national surveys used web, mail or mixed modes, e.g.\ the WVS in the United States in 2017 and in Japan in 2019, or the EVS in Switzerland in 2017). All statistics use the survey weights provided by the WVS and the EVS.

The WVS contains two questions on subjective well-being:
\begin{itemize}
    \item \textit{Happiness}: ``Taking all things together, would you say you are: \textit{Very happy}; \textit{Quite happy}; \textit{Not very happy}; \textit{Not at all happy}.''
    \item \textit{Life satisfaction}: ``All things considered, how satisfied are you with your life as a whole these days?'', answered on a scale from 1 (\textit{Completely dissatisfied}) to 10 (\textit{Completely satisfied}).
\end{itemize}
Non-responses are rare (\nonresponseSatisfactionWVS\% for satisfaction and \nonresponseHappinessWVS\% for happiness on average) and uncorrelated with GDP per capita (correlation: \corNonresponseGDP); they are excluded.

\ifshowregion
\paragraph{Indicators of national well-being.}
Because there is no consensus on how to aggregate ordinal answers into a national indicator, and because the ranking of countries can depend on this choice \citep{bond_sad_2019}, I compute eight indicators for each country-year:
\begin{itemize}
    \item \textit{Happy}: share answering \textit{Quite} or \textit{Very happy};
    \item \textit{Very Happy}: share answering \textit{Very happy};
    \item \textit{Very Unhappy}: share answering \textit{Not at all happy};
    \item \textit{V.~Happy -- V.~Unhappy}: difference between the two previous shares;
    \item \textit{Happiness (mean)}: mean happiness, with answers coded $+3$, $+1$, $-1$, $-3$;
    \item \textit{Satisfaction (mean)}: mean life satisfaction;
    \item \textit{Satisfied}: share answering 6 to 10 to the life-satisfaction question;
    \item \textit{Happy + Satisfied}: average of \textit{Happy} and \textit{Satisfied}, the index used by \citet{inglehart_genes_2000}.
\end{itemize}
For the life-satisfaction question, I also compute the shares answering 8--10, 9--10, 10, 1--4 and 1--2, which describe the tails of the distribution and are used for comparisons with Gallup. Finally, I compute the share of people whose satisfaction is below 60\% of the national mean, by analogy with relative poverty lines. The indicators are positively but imperfectly correlated: across country-years, the correlation between \textit{Happy} and \textit{Satisfied} is \corHappySatisfied, and between \textit{Very Happy} and \textit{Satisfaction (mean)} it is only \corVeryHappySatisfiedMean.
\fi

\subsection{Gallup World Poll and World Happiness Report}\label{subsec:gallup}

The Gallup World Poll interviews around 1,000 adults per country and year in more than 140 countries, by telephone where coverage is high (typically in high-income countries) and face-to-face elsewhere \citep{gallup_methodology_2022}. Its life-evaluation question is the Cantril ladder \citep{cantril_pattern_1965}: ``Please imagine a ladder, with steps numbered from 0 at the bottom to 10 at the top. The top of the ladder represents the best possible life for you and the bottom of the ladder represents the worst possible life for you. On which step of the ladder would you say you personally feel you stand at this time?''

I use the distribution of answers by country and wave, tabulated from the Gallup World Poll microdata (\nGallupCountryYears country-years, \nGallupCountries countries, 2006--2023), from which I compute the same satisfaction indicators as for the WVS: the mean and the shares answering 6--10, 8--10, 9--10, 10, 0--4 and 0--2. These distributions are unweighted. % Gallup's wave $w$ corresponds to year $w+2005$: with this mapping, three-year averages of the ladder computed from these distributions correlate at \corGallupWHR with the (weighted) three-year averages published in the World Happiness Report, against \corGallupWHRoldMapping with the mapping used in earlier versions of this project (year $w+2000$). 
To cover the most recent years, I also use the ladder means published in the World Happiness Report 2026 \citep{whr_2026}, which average Gallup data over three years up to 2023--2025.

\subsection{Income}\label{subsec:income}

My benchmark income measure is the log of GDP per capita at purchasing power parity (PPP, constant 2021 international dollars) from the World Development Indicators \citep{worldbank_wdi_2026}. As an alternative, I use GDP per capita at market exchange rates (constant 2015 US dollars). Missing values are filled in automatically, in three steps: (i) PPP series, which start in 1990, are extended backwards (and gaps filled) using the growth rate of real GDP per capita; (ii) for countries absent from the World Bank data (e.g. Taiwan, or Venezuela in recent years), I use the IMF World Economic Outlook \citep{imf_weo_2026}, whose values are in current dollars; I convert them into the constant dollars of the World Bank series by multiplying them by the ratio, in the same year, between the World Bank and the IMF values for the United States (a U.S. price deflator, since PPP and market-exchange-rate values coincide for the United States); (iii) remaining gaps are filled with the closest year available (Montenegro 1996). Territories without national accounts use the GDP of the closest entity (e.g. the United Kingdom for Northern Ireland). A robustness check excludes all imputed values.

\ifshowregion
Region is a categorical variable, which can capture non-linear patterns that a single continuous variable cannot. To compare income and region on an equal footing, I also use categorical income indicators: income sextiles and income clusters obtained by $k$-means clustering of log GDP per capita, with $k = 5$, 6 or 7. The clusters let the data choose the thresholds between income groups and allow for non-linear relationships, at the cost of more parameters than the log-linear specification. Sextiles and clusters are recomputed in each estimation sample: for instance, when the analysis is restricted to wave 3, the income groups are defined from the country-years surveyed in wave 3 only.
\fi

\ifshowregion
\subsection{Regions}\label{subsec:regions}

My main classification groups countries into the five UN regional groups (Figure~\ref{fig:regions}): \textit{Africa}, \textit{Asia} (the Asia-Pacific group, which includes the Middle East and Central Asia), \textit{Eastern Europe} (the former Eastern Bloc, excluding Central Asia), \textit{Latin America} (including the Caribbean), and \textit{Western} countries (the ``Western European and Others'' group: Western Europe, North America, Australia, New Zealand, and Turkey). It departs from the UN groups for two countries only: Israel, a member of the Western group since 2000 because it was denied access to the Asia-Pacific group, is classified in Asia, and Cyprus, a member of the Asia-Pacific group, is classified with Western countries. 
Territories that are not UN members (e.g.\ Taiwan, Hong Kong, Puerto Rico, Kosovo) are assigned to the group of their neighbors. For robustness, I use five alternative classifications: the exact UN regional groups (with Israel in the Western group and Cyprus in Asia); a six-region variant separating the Middle East and grouping Central Asia with the former Eastern Bloc; the seven World Bank regions; the five continents; and the UN geoscheme sub-regions (about twenty groups).

\begin{figure}[htbp]
    \centering
    \caption{Main classification of countries into five world regions}\label{fig:regions}
    \includegraphics[width=.9\textwidth]{../figures/region_groupings.png}
\end{figure}
\fi

% Shared section: the Fabre (2025) survey experiment
\subsection{A survey experiment on question wording and scale}\label{subsec:fabre}

Gallup and the WVS differ in many respects. Gallup asks the Cantril ladder on a 0--10 scale, the WVS life satisfaction on a 1--10 scale. Gallup relies on a centralized, standardized design, with about 1,000 respondents per country every year, interviewed by telephone in high-income countries and face-to-face elsewhere; the WVS is fielded every five years or so by national teams, mostly face-to-face, with sampling designs, modes and survey years that vary across countries. These differences fall into two categories that could explain why the two datasets disagree: the question (wording and scale) and everything else---sampling frames, interview modes, questionnaire context, translations, and survey years. To separate the two, I embedded a randomized experiment in an original survey \citep{fabre_public_2025}. The survey was conducted online between April 15 and July 3, 2025, by Bilendi (or Kantar in Saudi Arabia) on \nFabre respondents in \nFabreCountries high-income countries: the United States (3,000), Japan (2,000), Saudi Arabia (1,000), Germany (1,048), the United Kingdom (826), France (798), Italy (756), Spain (603), Poland (500) and Switzerland (469).\footnote{The survey also covered Russia, where the well-being question was censored.} Samples were stratified with quotas on gender, age, income, education, region and urbanicity,\footnote{Quota variables also included race in the U.S. and citizenship in Saudi Arabia.} and results are reweighted to match population frequencies along these dimensions (with weights trimmed between 0.25 and 4). %Respondents failing an attention test or completing the questionnaire in less than 6 minutes were excluded. The project was approved by the CIRED institutional review board (IRB-CIRED-2025-2) and pre-registered at the Open Science Framework (\href{https://osf.io/7mzn4}{osf.io/7mzn4}).

Toward the end of the questionnaire, after questions on global policies, each respondent was randomly assigned to one of four versions of the life-evaluation question:
\begin{itemize}
    \item \textit{Ladder 0--10} (Gallup's original): the Cantril ladder question quoted in Section~\ref{subsec:gallup}, with steps numbered from 0 (\textit{Worst possible}) to 10 (\textit{Best possible});
    \item \textit{Ladder 1--10}: same wording, with steps numbered from 1 to 10;
    \item \textit{Satisfaction 0--10}: the WVS question quoted in Section~\ref{subsec:wvs}, with answers from 0 (\textit{Completely dissatisfied}) to 10 (\textit{Completely satisfied});
    \item \textit{Satisfaction 1--10} (WVS's original): same wording, with answers from 1 to 10.
\end{itemize}
This $2 \times 2$ design identifies the effect of wording (ladder vs.\ satisfaction) and scale (0--10 vs.\ 1--10), and their interaction, holding the sample, mode, context and time constant. Within each country, the four branches are balanced (each receives between around 25\% % 21\% and 30\% 
of respondents).


\section{Income and region as predictors of national well-being}\label{sec:region}
% Shared section: income, region and national well-being (method and results)
\subsection{Empirical strategy}\label{subsec:method_region}

For each well-being indicator and each income indicator, I estimate three linear regressions across country-years $i$:
\begin{align}
    well\text{-}being_i &= \alpha_1 + \beta_1\, income_i + u_i \label{eq:income}\\
    well\text{-}being_i &= \alpha_2 + \gamma_2\, region_i + e_i \label{eq:region}\\
    well\text{-}being_i &= \alpha_3 + \beta_3\, income_i + \gamma_3\, region_i + \varepsilon_i \label{eq:both}
\end{align}
where $region_i$ denotes the vector of region dummies (so that $\gamma_2$ and $\gamma_3$ are vectors of coefficients). Denoting $R^2_1$, $R^2_2$ and $R^2_3$ their coefficients of determination, the share of the variance explained by income alone is $R^2_1$. To compare income and region, I decompose $R^2_3$ between the two regressors following \citet{lindeman_introduction_1980} (LMG), i.e.\ using the Shapley value \citep{gromping_estimators_2007}. With two (groups of) regressors, the LMG contribution of income is the average of its contribution when entered first ($R^2_1$) and when entered second ($R^2_3 - R^2_2$). The \textit{share of explained variance due to income} is therefore
\begin{equation}
    s = \frac{R^2_1 + (R^2_3 - R^2_2)}{2\, R^2_3}, \label{eq:share}
\end{equation}
and region is the better predictor whenever $s < 50\%$.

One feature of this comparison could favor region: region has four degrees of freedom while log GDP has one. I address this concern in three ways. First, I use categorical income variables (sextiles and clusters) with as many or more categories as regions. Second, I report the share of explained variance due to income computed with adjusted $R^2$, which penalizes additional regressors. I keep the unadjusted $R^2$ as the benchmark because it is the statistic used in the literature, and because the adjustment only reinforces the conclusions: with more than 300 observations, the adjustment is small, and with adjusted $R^2$ region is the better predictor for every indicator and income measure of Table~\ref{tab:share_income}. Third, and more directly relevant to the question of \textit{predictive} capacity, I compute leave-one-country-out cross-validated $R^2$: each country's observations are predicted by a model estimated on all other countries, and the cross-validated $R^2$ measures the share of the variance of well-being that the model correctly predicts for countries it was not estimated on. Unlike the in-sample $R^2$, it decreases when a model has superfluous parameters, and it is negative when a model predicts worse than the sample mean. Besides, as results could depend on how countries are grouped, I report results for five alternative classifications. Other robustness checks include each WVS wave separately, excluding the EVS surveys, weighting countries by their population, keeping only the last observation of each country, excluding the pandemic years 2020--2021, excluding imputed GDP values, using the previous vintage of GDP data (PPP in constant 2017 dollars), and excluding Latin America and Eastern Europe.

\subsection{Income explains a small share of national well-being in the WVS}\label{subsec:r2_income}

Figure~\ref{fig:satisfaction_gdp} plots mean life satisfaction against GDP per capita for all WVS country-years. The relationship is positive, but the dispersion at any given income level is large: countries with a GDP per capita around 10,000 dollars span most of the range of observed life satisfaction, from the former Soviet republics at the bottom to Latin American countries at the top. Log GDP per capita explains \RtwoSatisfiedMeanPPP\% of the variance of mean satisfaction.
% TODO refer to te Appendix Figures

\begin{figure}[htbp]
    \centering
    \caption{Mean life satisfaction vs.\ GDP per capita (PPP), WVS/EVS 1981--2023.}\label{fig:satisfaction_gdp}
    \includegraphics[width=.95\textwidth]{../figures/main/wvs_satisfied_mean_vs_gdp_ppp.pdf}
    % \begin{minipage}{.95\textwidth}\footnotesize \textit{Note:} Each point is a country-year (label: ISO code and last two digits of the year). The legend gives the $R^2$ of the regression on log GDP per capita.\end{minipage}
\end{figure}

Table~\ref{tab:r2_income} generalizes this result to all indicators and income measures. On average over the eight indicators, log GDP per capita at PPP explains \meanRtwoIncomePPP\% of the variance, and the best-performing income measure (\bestIncome) explains \meanRtwoIncomeBest\%. The indicators are positively but imperfectly correlated (Appendix Table~\ref{tab:correlation}). Satisfaction-based indicators are better explained than happiness-based ones: log GDP explains \RtwoHappinessMeanPPP\% of the variance of mean happiness and essentially none of the share of \textit{Very happy} people, whose slope with respect to log GDP is statistically indistinguishable from zero (slope \slopeVeryHappy, $p \pVeryHappy$; Appendix Table~\ref{tab:slopes}). This confirms the observation of \citet{deaton_income_2008} and \citet{inglehart_development_2008} that income is more closely related to life evaluation than to happiness. By contrast, region alone explains \meanRtwoRegion\% of the variance on average, and scores high even on happiness-based variables (last column).

\begin{table}[htbp]
    \centering
    \caption{Share of the variance of national well-being explained by income alone ($R^2$), WVS/EVS 1981--2023.}\label{tab:r2_income}
    \begin{adjustbox}{max width=\textwidth}
    \input{../tables/main/r2_income.tex}
    \end{adjustbox}
    \begin{minipage}{\textwidth}\footnotesize \textit{Note:} $R^2$ of the regression of each well-being indicator (rows) on each income indicator (columns), over \nWVSCountryYears WVS/EVS country-years (a few with missing GDP are excluded). The last column gives the $R^2$ of the regression on region dummies. The replication files include an Excel version of this table with the $R^2$ of all region classifications (\texttt{tables/main/r2\_income\_all\_regions.xlsx}).\end{minipage}
\end{table}

\subsection{Region predicts national well-being better than income in the WVS}\label{subsec:region_wvs}

Table~\ref{tab:share_income} reports the share of explained variance due to income, $s$, for each indicator and income measure. It is below 50\% in all cases. On average, the best income measure accounts for \meanShareIncomeBest\% of the variance explained jointly by income and region. The first split of a regression tree using log GDP and region as candidate variables is on region for all %\treeRegionFirst out of 8 
national well-being indicators.

% TODO shouldn't we have "." at the end of captions?
\begin{table}[htbp]
    \centering
    \caption{Share of explained variance due to income rather than region (LMG), WVS/EVS 1981--2023}\label{tab:share_income}
    \begin{adjustbox}{max width=\textwidth}
    \input{../tables/main/share_income.tex}
    \end{adjustbox}
    \begin{minipage}{\textwidth}\footnotesize \textit{Note:} Share of the $R^2$ of Equation~\eqref{eq:both} attributed to income by the LMG (Shapley) decomposition, Equation~\eqref{eq:share}. Values below 0.5 indicate that region is the better predictor.\end{minipage}
\end{table} % TODO put cells above .5 in bold, and say it's the case in note

The result is not an artifact of the larger number of parameters of the region model. Out of sample, region predicts better than income for every indicator (Table~\ref{tab:cv}): for mean satisfaction, the leave-one-country-out $R^2$ is \cvRegionSatisfiedMean with region against \cvIncomeSatisfiedMean with log GDP. For several happiness indicators, income has no out-of-sample predictive power at all (negative cross-validated $R^2$).

\begin{table}[htbp]
    \centering
    \caption{Out-of-sample predictive capacity: leave-one-country-out cross-validated $R^2$}\label{tab:cv}
    \begin{adjustbox}{max width=\textwidth}
    \input{../tables/main/cross_validated_r2.tex}
    \end{adjustbox}
    \begin{minipage}{\textwidth}\footnotesize \textit{Note:} Each country's observations are predicted with a model estimated on all other countries. Cross-validated $R^2 = 1 - \sum (y - \hat y_{\text{out-of-sample}})^2 / \sum (y - \bar y)^2$; it can be negative when a model predicts worse than the overall mean.\end{minipage}
\end{table} % TODO put in bold the higher cell of each row among columns log GDP, best income, and region, and say you do so in the table note

Table~\ref{tab:robustness} summarizes \nSpecsWVS estimates (indicator $\times$ income measure $\times$ specification) % TODO clarify what "specification" is, and replace by "dataset" if this is what it means
using the WVS. Region is the better predictor in \shareRegionBetterWVS\% of them (\shareRegionBetterWVSbest\% when restricting to the best income measure). The result is robust. It holds in each wave except wave 4 (1999--2004, 40 countries), where income and region perform similarly; when countries are weighted by population; without the EVS surveys, or when keeping one observation per country in case there is both an EVS and WVS in one year; without the pandemic years; without imputed GDP values; with the previous vintage of GDP data; and with the exact UN groups, the six-region or the UN sub-region classifications.

\begin{table}[htbp]
    \centering
    \caption{Robustness: income vs.\ region across specifications and datasets}\label{tab:robustness}
    \begin{adjustbox}{max width=\textwidth}
    \input{../tables/main/robustness.tex}
    \end{adjustbox}
    \begin{minipage}{\textwidth}\footnotesize \textit{Note:} Averages over well-being indicators (8 for the WVS; 7 satisfaction indicators for Gallup; mean ladder for the WHR). ``Best'' refers to the income measure with the highest average $R^2$ in the main WVS specification (\bestIncome). ``Region better'' is the share of indicator $\times$ income-measure combinations for which $s < 0.5$.\end{minipage}
\end{table} % TODO put in bold the cells >0.5 in the last column

\subsection{Where the region advantage comes from}\label{subsec:region_limits}

Three results qualify this conclusion. First, the advantage of region is driven by two well-known ``anomalies'': Latin America, happier than its income would predict \citep{graham_paradox_2009,rojas_handbook_2016}, and the former Eastern Bloc, unhappier \citep{guriev_unhappiness_2009}. Excluding these two regions, region is the better predictor in only \shareRegionBetterNoLatamEE\% of cases. Second, the result depends on the granularity and design of the classification: with continents or World Bank regions, which do not isolate Latin America or the former Eastern Bloc, region is better in \shareRegionBetterContinents\% and \shareRegionBetterWB\% of cases, respectively. Third, and most importantly, the conclusion does not carry over to Gallup data. In Gallup's latest observation for each of \nGallupCountries countries, log GDP explains \RtwoGallupMeanPPP\% of the variance of the mean ladder, against \RtwoGallupRegion\% for region, and income accounts for \shareIncomeGallupMean\% of the explained variance (Appendix Table~\ref{tab:share_gallup}). Over all Gallup specifications, region is the better predictor in only \shareRegionBetterGallup\% of cases, mostly for indicators of the upper tail of the distribution (shares answering 9--10 or 10), which are almost unrelated to income. With the latest WHR data (2023--2025), log GDP explains \RtwoWHRPPP\% of the variance of the ladder against \RtwoWHRRegion\% for region.

The relation between GDP and well-being is thus much weaker in the WVS than in Gallup. The same holds when restricting both datasets to the \nCommon country-years observed by both surveys: the $R^2$ of mean life evaluation on log GDP is \RtwoCommonGallup in Gallup and \RtwoCommonWVS in the WVS. Whether ``GDP is a poor predictor of national well-being'' therefore depends on which of the two datasets one trusts.

\subsection{Other correlates and within-country evidence}\label{subsec:correlates}

What do regions capture? Appendix Table~\ref{tab:correlates} decomposes the variance of mean satisfaction among income, region and one additional country-level variable from the WVS. The average perceived freedom of choice and control over one's life (``how much freedom of choice and control you feel you have over the way your life turns out'') alone explains \RtwoFreedomSatisfaction\% of the variance of mean satisfaction; in a Shapley decomposition with income and region, it accounts for \shapleyFreedom percentage points, against \shapleyFreedomRegion for region and \shapleyFreedomIncome for income. This variable is itself a subjective evaluation, so this correlation should not be read causally, but it suggests that non-material dimensions of life are central to regional differences \citep{inglehart_development_2008}. Tolerance of homosexuality (\RtwoTolerance\%), the importance of God (\RtwoGod\%) and GDP growth over the previous five years (\RtwoGrowth\%) explain much less.

Finally, the cross-sectional relation should not be confused with the relation over time. Within countries (country fixed effects, \nWithin observations from countries surveyed at least twice), mean satisfaction increases with log GDP per capita, with a slope of \withinSlopeSatisfiedMean per unit of $\log_{10}$ GDP, larger than the cross-sectional slope (\crossSlopeSatisfiedMean), and even the share of \textit{Very happy} people increases (slope \withinSlopeVeryHappy, $p \pWithinVeryHappy$). Income growth explains \withinRtwoSatisfiedMean of the within-country variance (Appendix Table~\ref{tab:within}). These within-country associations are consistent with \citet{stevenson_economic_2008} and \citet{sacks_new_2012} rather than with a strict version of the Easterlin paradox \citep{easterlin_does_1974,easterlin_happiness_2010}, although they may reflect common shocks (e.g.\ transition recessions) rather than the effect of income per se.

\subsection{What are the happiest countries?}\label{subsec:happiest}

Rankings of countries depend on the indicator. Taking, for each indicator and each wave (and all waves combined), the happiest country-year, \happiestFirst appears most often (\happiestFirstN times out of \nHappiestCells), followed by \happiestSecond (\happiestSecondN) and \happiestThird (\happiestThirdN) (Appendix Table~\ref{tab:happiest}). The happiest country-years are Western (\happiestRegionWestern), Latin American (\happiestRegionLatinAmerica), Asian (\happiestRegionAsia) or African (\happiestRegionAfrica). The richest countries are thus not necessarily the happiest: the link between income and well-being seems to be that rich countries rarely experience low well-being, rather than that they reach the highest levels.


\section{Why do Gallup and the WVS disagree? Wording vs.\ samples}\label{sec:discrepancy}
% Shared section: why do Gallup and the WVS disagree? Wording vs. sampling
\subsection{Experimental effects of wording and scale}\label{subsec:experiment}

Table~\ref{tab:fabre_regressions} reports regressions of the answer on the randomized treatments, with country fixed effects, survey weights and heteroskedasticity-robust standard errors. The ladder wording shifts answers by \effectLadder points on average (s.e.\ \seLadder), i.e.\ respondents rate their lives lower on the ``best possible life'' ladder than when asked how satisfied they are. This is in line with evidence that the two questions elicit different considerations: in open-ended answers, the ladder evokes power and wealth more than life satisfaction does \citep{nilsson_cantril_2024}. The scale effect is smaller: answers on the 1--10 scale are \effectScale points higher than on the 0--10 scale (s.e.\ \seScale), which, once the 1--10 answers are linearly stretched to 0--10, amounts to \effectScaleZeroTen points. Wording and scale do not interact. Measured as the share of respondents answering 6 or more, the ladder changes the share ``satisfied'' by \effectLadderSatisfied percentage points.

\begin{table}[htbp]
    \centering
    \caption{Effects of question wording and scale on life evaluation, Fabre (2025) survey experiment}\label{tab:fabre_regressions}
    \begin{adjustbox}{max width=\textwidth}
    \input{../tables/main/fabre_regressions.tex}
    \end{adjustbox}
    \begin{minipage}{\textwidth}\footnotesize \textit{Note:} OLS with country fixed effects, survey weights (normalized within country) and heteroskedasticity-robust standard errors in parentheses. \textit{Ladder wording}: Cantril ladder (vs.\ WVS life satisfaction). \textit{Scale 1--10}: answers from 1 to 10 (vs.\ 0 to 10). In columns 3 and 5, 1--10 answers are linearly rescaled to 0--10. Column 5 interacts the wording with the respondent's income quartile within their country. * $p<.05$, ** $p<.01$, *** $p<.001$.\end{minipage}
\end{table}

Two further results matter for the interpretation of cross-country comparisons. First, the ladder does not steepen the income gradient \textit{within} countries: the interaction between ladder wording and the income quartile is \interactionLadderIncome ($p \pInteractionLadderIncome$), against a gradient of \gradientIncomeSatisfaction points per quartile (column 5). The scale does not change this conclusion: in a model interacting income with both wording and scale (column 6), the interaction between income and wording is \interactionLadderIncomeFull ($p \pInteractionLadderIncomeFull$) and the triple interaction is \tripleInteraction ($p \pTripleInteraction$). So, at least among individuals in high-income countries, the ladder is not more ``about money'' than the satisfaction question. Second, the wording effect differs markedly across countries (Figure~\ref{fig:fabre_variants}). It ranges from \wordingMin points (Spain) to \wordingMax (Japan), where the ladder yields \textit{higher} answers than the satisfaction question, and the hypothesis of a common effect across countries is strongly rejected ($\chi^2(9) = \chiHeterogeneityWording$, $p \pHeterogeneityWording$). Question wording therefore affects not only the level but also the ranking of countries.

\begin{figure}[htbp]
    \centering
    \caption{Mean answer by question variant and country, Fabre (2025)}\label{fig:fabre_variants}
    \includegraphics[width=.85\textwidth]{../figures/main/fabre_variants_by_country.pdf}
    \begin{minipage}{.85\textwidth}\footnotesize \textit{Note:} Weighted mean answer on the native scale of each variant (0--10 or 1--10).\end{minipage}
\end{figure}

\subsection{Decomposing the Gallup--WVS gap in ten countries}\label{subsec:decomposition}

For each country $c$ of the experiment, I define the observed gap between Gallup and the WVS as $D_c = G_c - W_c$, where $W_c$ is the mean satisfaction in the latest WVS or EVS survey of the country (from 2017--2022, except for Saudi Arabia, 2003) and $G_c$ the mean ladder in Gallup the same year (or the closest year, at most three years apart; \nSameYear countries have an exact match). Comparing the two datasets in the same year removes the influence of changes in well-being over time. The experiment measures, within the same sample, the \textit{question effect} $Q_c = F_c(\text{ladder 0--10}) - F_c(\text{satisfaction 1--10})$, where $F_c(v)$ is the mean answer to variant $v$ in country $c$. The residual $R_c = D_c - Q_c$ is the part of the gap due to everything that differs between the two datasets except the question: sampling frames, interview modes, questionnaire context, translations, and survey-specific time effects. This decomposition assumes that the question effect measured online in 2025 applies to the Gallup and WVS samples of earlier years.

Table~\ref{tab:decomposition_country} and Figure~\ref{fig:decomposition} present the decomposition for each country, and Table~\ref{tab:decomposition_summary} summarizes it with 95\% confidence intervals from a bootstrap that resamples respondents in the three surveys. On average, the gap between Gallup and WVS answers is $D = \decompMeanGap$ points. The question effect is even larger in absolute value: $Q = \decompMeanQuestion$ points. Wording (and scale) thus fully account for the fact that Gallup yields lower levels than the WVS in high-income countries, and the residual is small and not significantly different from zero ($R = \decompMeanResidual$, 95\% CI [\decompLowMeanResidual; \decompHighMeanResidual]). Since the two components move in opposite directions, another way to summarize the level difference is to compare the size of the two movements: the question accounts for $|\bar Q|/(|\bar Q|+|\bar R|) = \decompShareAbsQuestion$ (95\% CI [\decompLowShareAbsQuestion; \decompHighShareAbsQuestion]) of the gross difference in levels between Gallup and the WVS, and for \decompShareAbsQuestionCountry on average when this share is computed country by country. By this metric, wording is the main driver of the difference in \textit{levels}.

\begin{table}[htbp]
    \centering
    \caption{Decomposition of the Gallup--WVS gap by country}\label{tab:decomposition_country}
    \begin{adjustbox}{max width=\textwidth}
    \input{../tables/main/decomposition_by_country.tex}
    \end{adjustbox}
    \begin{minipage}{\textwidth}\footnotesize \textit{Note:} Mean answers on native scales (Gallup: 0--10; WVS/EVS: 1--10). The Gallup year is given in parentheses when it differs from the WVS year. Survey: mode: source of the WVS/EVS data and main interview mode(s) (CAPI: computer-assisted personal interviews).\end{minipage}
\end{table}

\begin{figure}[htbp]
    \centering
    \caption{Observed Gallup--WVS gap, question effect and residual by country}\label{fig:decomposition}
    \includegraphics[width=.9\textwidth]{../figures/main/decomposition_by_country.pdf}
\end{figure}

What matters for correlations with GDP and for country rankings, however, is not the level of the gap but its variation across countries, which the previous metrics do not capture: a question effect that shifted all countries by the same amount would score high on them while leaving correlations and rankings unchanged. Here, the question effect explains little: the share of the cross-country variance of the gap attributable to the question, $\mathrm{cov}(D,Q)/\mathrm{var}(D)$, is \decompVarShareQuestion (95\% CI [\decompLowVarShareQuestion; \decompHighVarShareQuestion]), and the residual accounts for the rest. Predicting a country's Gallup score by adding its experimentally measured question effect to its WVS score does worse (root mean squared error: \decompRmsePredictionGallup) than simply adding the average gap (\decompRmseNaive). The conclusion is the same when excluding Saudi Arabia, the only country whose latest WVS survey is not recent (variance share: \varShareQuestionRecent), and for the share of respondents answering 6 or more (\decompVarShareQuestionSat). For instance, Japan is the only country with a positive question effect, but the gap is negative there; and in the United States, the gap is nil although the question effect is as large as in other countries.

\begin{table}[htbp]
    \centering
    \caption{Decomposition of the Gallup--WVS gap: summary}\label{tab:decomposition_summary}
    \begin{adjustbox}{max width=\textwidth}
    \input{../tables/main/decomposition_summary.tex}
    \end{adjustbox}
    \begin{minipage}{\textwidth}\footnotesize \textit{Note:} \nDecompCountries countries. 95\% confidence intervals from 1,000 bootstrap replications resampling respondents within country $\times$ variant (Fabre), within country (WVS), and answers within country (Gallup, multinomial).\end{minipage}
\end{table}

Two direct comparisons of samples asked the same question corroborate the importance of the residual. First, Gallup respondents (World Happiness Report, 2023--2025) rate their lives \meanSampleEffectGallup points higher on the 0--10 ladder than Fabre (2025) online respondents asked the same question, with differences ranging from \minSampleEffectGallup to \maxSampleEffectGallup across countries. Second, WVS respondents report \meanSampleEffectWVS points higher satisfaction than online respondents asked the WVS question (this comparison also includes changes over time, as WVS surveys are older). These sample (and mode) effects are of the same order of magnitude as the question effect, and they vary across countries. Even in the United States, where the 2017 WVS was itself conducted online, WVS respondents report \sampleEffectUSWVS points higher satisfaction than Fabre (2025) respondents.

\subsection{Income gradients in the ten countries: past vs.\ new data}\label{subsec:ten_countries}

The decomposition above concerns levels. The question that motivates this paper, however, is whether the two datasets disagree on the \textit{relationship} between income and well-being because of their questions. The ten countries of the experiment provide a direct test (Table~\ref{tab:ten_countries}). In the past data, the income gradient is much steeper in Gallup than in the WVS: across the ten countries, the slope of mean life evaluation on $\log_{10}$ GDP per capita is \slopeTenGallup in Gallup and \slopeTenWVS in the WVS, a difference of \SlopeDiffPast (95\% bootstrap CI [\SlopeDiffPastLow; \SlopeDiffPastHigh]), and log GDP explains \RtwoTenGallup of the cross-country variance in Gallup against \RtwoTenWVS in the WVS. In the new data, where the same sample answers either question, the gradient is similar with both wordings: \slopeTenFabreLadder with the ladder and \slopeTenFabreSatisfaction with the satisfaction question, a difference of \SlopeDiffFabre (95\% CI [\SlopeDiffFabreLow; \SlopeDiffFabreHigh]); the two other variants yield intermediate slopes. Both original questions thus yield a gradient close to Gallup's. The effect of wording on the gradient is imprecisely estimated, but it is significantly smaller than the difference between Gallup and the WVS: the part of this difference that is not due to the question is \SlopeDiffinDiff (95\% CI [\SlopeDiffinDiffLow; \SlopeDiffinDiffHigh]). The flat gradient of the WVS is thus mostly specific to its samples, not to its question. This is consistent with the individual-level evidence that the ladder does not steepen the income gradient within countries (Section~\ref{subsec:experiment}). Two caveats apply. First, the bootstrap confidence intervals reflect the sampling of respondents and treat the ten countries as given; with ten countries, the slopes themselves are imprecise (Table~\ref{tab:ten_countries} reports their standard errors across countries). Second, the $R^2$ differ more than the slopes, because the satisfaction question produces more dispersion around the income gradient (notably in Japan), and the decompositions between income and region are uninformative here, since Poland is the only Eastern European country.

\begin{table}[htbp]
    \centering
    \caption{Income gradient of national well-being in the ten countries of the experiment: past vs.\ new data}\label{tab:ten_countries}
    \begin{adjustbox}{max width=\textwidth}
    \input{../tables/main/ten_countries.tex}
    \end{adjustbox}
    \begin{minipage}{\textwidth}\footnotesize \textit{Note:} OLS across the \nDecompCountries countries of the experiment. Gallup and WVS/EVS: latest WVS/EVS survey and Gallup the same year, GDP of that year. Fabre (2025): mean answer by variant, GDP of 2025. Standard errors robust to heteroskedasticity (HC1). $R^2$ income, $R^2$ region and share due to income: see Section~\ref{subsec:method_region} (regions: Western, Asia, Eastern Europe).\end{minipage}
\end{table}

Two further comparisons of samples asked the same question confirm that sample differences are large. First, the EVS and the WVS surveyed \nEVSvsWVS countries separately during 2017--2022, with the same satisfaction question; their mean satisfaction differs by \meanAbsDiffEVSWVS points on average (up to \maxAbsDiffEVSWVS), a magnitude comparable to the cross-country variation of the Gallup--WVS gap. Second, as noted above, Gallup and WVS respondents report markedly higher well-being than online respondents asked the same question.

\subsection{The global discrepancy}\label{subsec:global_gap}

The ten countries of the experiment are all high-income, whereas the Gallup--WVS discrepancy matters most for the global relation between GDP and well-being. Restricting both datasets to the \nCommon country-years (\nCommonCountries countries) observed by both surveys the same year, log GDP explains \RtwoCommonGallup of the variance of Gallup's mean ladder but only \RtwoCommonWVS of WVS mean satisfaction, and income accounts for \shareIncomeCommonGallup\% of the variance explained jointly with region in Gallup against \shareIncomeCommonWVS\% in the WVS. The gap between the two datasets increases strongly with GDP (Figure~\ref{fig:gap_gdp}): the slope is \slopeGapGDP per unit of $\log_{10}$ GDP (s.e.\ \seSlopeGapGDP, clustered by country), with WVS answers exceeding Gallup answers by around two points in low-income countries, against less than half a point in high-income countries.

\begin{figure}[htbp]
    \centering
    \caption{Gallup ladder minus WVS satisfaction in the same country-year, by GDP per capita}\label{fig:gap_gdp}
    \includegraphics[width=.9\textwidth]{../figures/main/gap_gallup_wvs_vs_gdp.pdf}
    \begin{minipage}{.9\textwidth}\footnotesize \textit{Note:} Mean answers on native scales (Gallup: 0--10; WVS: 1--10), \nCommon country-years observed by both surveys. Line: OLS fit with 95\% confidence band.\end{minipage}
\end{figure}

Could the question explain this pattern? Within the experiment, the question effect is not significantly related to GDP (slope \slopeQuestionGDP, s.e.\ \seSlopeQuestionGDP), which is unsurprising given the narrow range of GDP among the ten countries. Even taking this point estimate at face value and extrapolating it to all countries, adjusting WVS answers for the predicted question effect would raise the $R^2$ of WVS satisfaction on log GDP only from \RtwoCommonWVS to \RtwoAdjustedWVS, far from Gallup's \RtwoCommonGallup. The experiment cannot rule out that the wording effect is much larger in low-income countries---for instance if the ``best possible life'' evokes material comparisons more strongly where material deprivation is widespread---but nothing in the high-income evidence points in that direction: the ladder does not steepen the individual income gradient (Section~\ref{subsec:experiment}), and the country-level question effects are unrelated to the observed gaps.


\subsection{Which dataset has more representative samples?}\label{subsec:data_quality}

Since the discrepancy stems from samples rather than questions, which dataset is closer to the truth? The methodological documentation offers arguments on both sides. Gallup applies a centralized, standardized design in all countries: probability samples of about 1,000 adults per year, drawn by random-digit dialing where telephone coverage is high and by multistage area sampling elsewhere, with random selection of the respondent within the household, post-stratification weights on gender, age and, where possible, education, and documented coverage exclusions \citep{gallup_methodology_2023}. But its samples are small; its face-to-face samples select households by random route, with replacement of non-responding households; response rates are not published; interviews switched from face-to-face to telephone during the pandemic (2020--2021); and in 2023, a fifth of the interviews in 27 high-income countries came from web panels, most of them opt-in \citep{gallup_methodology_2023}. \citet{blanchflower_gallup_2024} also questions the reliability of Gallup estimates for subgroups within countries. The WVS requires probability samples of at least 1,200 respondents but allows quota selection at the last stage, and, being fielded by national teams, its sampling designs, modes and survey years vary across countries \citep{wvs_rules_wave8}. Automated tests suggesting duplicated interviews in international survey programs \citep{kuriakose_dont_2016} have been disputed \citep{pew_data_2016} and cannot be applied to Gallup, whose microdata are not public.

I test the representativeness of WVS/EVS samples directly (Table~\ref{tab:representativeness}). Compared with population benchmarks \citep{barro_new_2013,worldbank_wdi_2026}, WVS samples over-represent tertiary-educated adults (aged 25 and over) by a median factor of \ratioTertiaryLow in countries with a GDP per capita below 10,000 dollars, against \ratioTertiaryHigh in countries above 30,000 dollars (correlation of the log ratio with log GDP: \corTertiaryGDP, \nCompositionSurveys surveys). They also over-represent urban residents in poorer countries (median ratio \ratioUrbanLow vs.\ \ratioUrbanHigh, \nUrbanSurveys recent surveys) and under-represent people aged 65 or more (\ratioOldLow vs.\ \ratioOldHigh). The internal criterion of \citet{kohler_surveys_2007}---women should make up half of married or cohabiting respondents---is violated at the 5\% level in \shareGenderFlag\% of surveys (\shareGenderFlagLow\% in countries below 10,000 dollars and \shareGenderFlagHigh\% above 30,000 dollars), with a median deviation of \medianGenderDeviation percentage points. WVS samples are thus less representative in poorer countries, which is where their gap with Gallup is largest. However, these observable imbalances do not explain the gap. Controlling for log GDP, the over-representation of the tertiary-educated does not predict the Gallup--WVS gap (coefficient \coefTertiaryGap, s.e.\ \seTertiaryGap, $p \pTertiaryGap$, \nGapComposition surveys), and reweighting each WVS sample to the benchmark share of tertiary-educated adults barely changes its income gradient: in the country-years also covered by Gallup, the $R^2$ of mean satisfaction on log GDP goes from \RtwoCommonWVSraw to \RtwoCommonWVSeduAdj, against \RtwoCommonGallupSub for Gallup.

\begin{table}[htbp]
    \centering
    \caption{Representativeness of WVS/EVS samples: sample shares relative to population benchmarks}\label{tab:representativeness}
    \begin{adjustbox}{max width=\textwidth}
    \input{../tables/main/sample_representativeness.tex}
    \end{adjustbox}
    \begin{minipage}{\textwidth}\footnotesize \textit{Note:} One observation per WVS or EVS survey. Ratios: weighted sample share divided by the population share, medians across surveys (a ratio above 1 indicates over-representation). Benchmarks (World Bank WDI): share of the population aged 25+ with at least short-cycle tertiary education (UNESCO), or with complete or incomplete tertiary education \citep{barro_new_2013} before 2011; urban population share (national definitions; sample: urban and peri-urban respondents, waves 4, 5 and 7 only); share aged 65+ among the population aged 15+ (sample: 65+ among respondents, usually 18+). Women among married: unweighted share of women among married or cohabiting respondents minus 0.5. Correlations with log GDP per capita (PPP) of the log ratios (or of the deviation).\end{minipage}
\end{table}

Overall, three arguments favor Gallup for cross-country comparisons: its standardized design and weighting, the lower representativeness of WVS samples in poorer countries, and the fact that an independent online sample reproduces Gallup's income gradient in high-income countries (Section~\ref{subsec:ten_countries}). But the discrepancy cannot be traced to the observable composition of WVS samples, which suggests that it stems from unobserved differences in who responds, how, and in which context; and Gallup's recent changes of interview modes call for caution when comparing years. (The Gallup distributions used in this paper are unweighted, but their three-year averages correlate at \corGallupWHR with the weighted averages published in the World Happiness Report.)

\section{Discussion}\label{sec:discussion}
\subsection{Income vs.\ region}
% Shared section: interpretation and limitations of the region vs. income analysis
Several caveats apply to the comparison between income and region. First, region is not a mechanism: it bundles culture, history, institutions, social ties, and possibly cross-cultural differences in how people use response scales \citep{chen_response_1995,diener_personality_2003,exton_comparing_2015}. If regional differences partly reflect response styles, then region predicts \textit{measured} well-being without necessarily predicting \textit{experienced} well-being. Existing evidence suggests that self-reported well-being carries information about true well-being \citep{sandvik_subjective_1993,kaiser_scientific_2022}, and studies exploiting migrants suggest that national differences partly reflect a cultural component of well-being rather than mere translation artifacts \citep{senik_french_2014,exton_comparing_2015}, but the present analysis cannot separate these interpretations. Anchoring vignettes \citep{kapteyn_life_2010} or measures of emotional well-being would help.

Second, the advantage of region rests largely on Latin America and the former Eastern Bloc. These are well-documented deviations from the income gradient, but they are not independent of the classification: grouping countries into continents or World Bank regions, which do not isolate these groups, weakens the result. The fairest summary is that GDP per capita leaves much of the cross-country variation unexplained, and that a few large regional deviations explain more of it than income does in the WVS.

Third, the analysis pools country-years that are not independent, and $R^2$ decompositions carry no information on causality. The within-country evidence shows that growth is associated with rising life satisfaction, so the weak cross-sectional relation does not imply that income is irrelevant for well-being.

Fourth, and most importantly, the WVS and Gallup disagree on the strength of the income gradient. In Gallup, income predicts life evaluations about as well as region, or better.

Section~\ref{sec:discrepancy} showed that this discrepancy comes mostly from differences between samples in the broad sense, not from the question asked: asked to the same sample, the ladder and the satisfaction question yield similar income gradients.
\subsection{Wording vs.\ samples}
% Shared section: interpretation and limitations of the wording vs. sampling decomposition
The experiment shows that the question matters for levels: the Cantril ladder yields lower answers than the WVS life-satisfaction question, and this difference fully accounts for the lower average level of Gallup data in high-income countries. But the question does not explain the cross-country pattern of the Gallup--WVS gap, which is what drives differences in correlations with GDP and in country rankings: asked to the same sample, both questions yield similar income gradients, close to Gallup's. The pattern is explained by the residual, i.e.\ by differences between samples in the broad sense: sampling frames, interview modes \citep{dolan_happy_2016,conti_survey_2011}, questionnaire context \citep{deaton_financial_2012,deaton_understanding_2016}, translations, and time.

Several limitations should be kept in mind. First, the experiment covers ten high-income countries, whereas the gap varies most across income levels. The question effect could be different in low- and middle-income countries, where both Gallup and the WVS mostly interview face-to-face; only an experiment in these countries can settle this. Second, the residual lumps together several factors that the design cannot separate. In particular, the question effect was measured with the survey's own translations, toward the end of a questionnaire on global policies, and online; it may differ from the effect that would be obtained with official Gallup and WVS translations, at the beginning of an interview, and face-to-face or by phone. Third, the decomposition assumes that the question effect measured in 2025 applies to the WVS/EVS and Gallup surveys of 2017--2022 (2003 for Saudi Arabia). Fourth, with ten countries, cross-country statistics are imprecise, as shown by the bootstrap confidence intervals. Finally, online panel respondents report markedly lower well-being than Gallup and WVS respondents asked the same question, a sample and mode effect that deserves attention in its own right as online panels become common in well-being research.


\section{Conclusion}\label{sec:conclusion}

GDP per capita is a weaker predictor of national well-being than is often assumed, but how weak depends on the data. In the WVS, it explains a small share of the cross-country variance of well-being, less than a country's world region, and it explains essentially nothing of the share of very happy people. In the Gallup World Poll, however, income predicts life evaluations as well as region does. A survey experiment in ten high-income countries shows that the question matters for levels---the Cantril ladder yields lower answers than the life-satisfaction question---but that it has little effect on the relationship between income and well-being: asked to the same sample, both questions yield an income gradient close to Gallup's, whereas the WVS gradient is much flatter. The discrepancy between the two datasets therefore reflects differences in samples, modes, and context.

These findings have two implications. For research, cross-country conclusions about the income--well-being relationship should be checked across datasets, and the choice between Gallup and the WVS should be based on the quality and comparability of their samples rather than on their questions, since wording effects mostly shift levels. Gallup's standardized annual methodology is an argument in its favor for cross-country comparisons, and in the ten high-income countries studied, an independent online sample reproduces Gallup's income gradient rather than the WVS's. However, the gap between datasets is largest in low-income countries, where neither the present experiment nor existing evidence can yet tell which survey is closer to the truth. Replicating the experiment face-to-face in low- and middle-income countries is the natural next step. For policy, the fact that rich countries rarely have low well-being, while many countries at intermediate income levels are as happy as the richest, suggests that growth is neither necessary nor sufficient for high well-being, and that the non-material determinants of well-being that differ across regions deserve more attention in policy evaluation \citep{stiglitz_report_2009,frijters_happy_2020}.

\clearpage
\section*{Declaration of generative AI and AI-assisted technologies in the writing process}
\todo{During the preparation of this work the author used Claude (Anthropic) to write analysis code and draft parts of the manuscript. After using this tool, the author reviewed and edited the content as needed and takes full responsibility for the content of the publication. (Adapt to actual use.)}

\bibliography{wellbeing}

\clearpage
\appendix
\section{Additional results}\label{app:additional}
% Shared appendix: additional results on income, region and national well-being
\renewcommand{\thetable}{A\arabic{table}}
\renewcommand{\thefigure}{A\arabic{figure}}
\setcounter{table}{0}
\setcounter{figure}{0}

\begin{table}[htbp]
    \centering
    \caption{Correlation between national well-being indicators (WVS country-years)}\label{tab:correlation}
    \begin{adjustbox}{max width=\textwidth}
    \input{../tables/main/correlation_indicators.tex}
    \end{adjustbox}
\end{table}

\begin{table}[htbp]
    \centering
    \caption{Slope of national well-being indicators with respect to $\log_{10}$ GDP per capita (PPP)}\label{tab:slopes}
    \begin{adjustbox}{max width=\textwidth}
    \input{../tables/main/slopes.tex}
    \end{adjustbox}
    \begin{minipage}{\textwidth}\footnotesize \textit{Note:} OLS across country-years; standard errors clustered by country. Gallup: 2006--2023.\end{minipage}
\end{table}

\begin{table}[htbp]
    \centering
    \caption{Share of explained variance due to income rather than region (LMG), Gallup, last observation per country}\label{tab:share_gallup}
    \begin{adjustbox}{max width=\textwidth}
    \input{../tables/main/share_income_gallup.tex}
    \end{adjustbox}
    \begin{minipage}{\textwidth}\footnotesize \textit{Note:} \nGallupCountries countries, last Gallup observation available (2006--2023). Same definitions as Table~\ref{tab:share_income}.\end{minipage}
\end{table} % TODO put cells above .5 in bold, and say it's the case in note

\begin{table}[htbp]
    \centering
    \caption{Other country-level correlates of national well-being: Shapley decomposition of $R^2$ (WVS)}\label{tab:correlates}
    \begin{adjustbox}{max width=\textwidth}
    \input{../tables/main/correlates.tex}
    \end{adjustbox}
    \begin{minipage}{\textwidth}\footnotesize \textit{Note:} Country-year averages of WVS variables: perceived freedom of choice and control (1--10), importance of God in one's life (1--10), justifiability of homosexuality (1--10), importance of living in a democracy (1--10, waves 5--7), and average annual growth of real GDP per capita over the previous five years. Shapley decomposition of the $R^2$ of the regression on log GDP per capita, region dummies and the other variable.\end{minipage}
\end{table} % TODO put in bold the highest cell in each row among Income, Region, Other.
% TODO add at the end of each table/figure note: \hfill(Back~to~Section~\ref{[section where it is cited]})

\begin{table}[htbp]
    \centering
    \caption{Within-country relation between national well-being and $\log_{10}$ GDP per capita (country fixed effects, WVS)}\label{tab:within}
    \begin{adjustbox}{max width=\textwidth}
    \input{../tables/main/within_country.tex}
    \end{adjustbox}
    \begin{minipage}{\textwidth}\footnotesize \textit{Note:} Countries surveyed at least twice. Standard errors clustered by country.\end{minipage}
\end{table} % TODO add the average R² over indicators in the table note

\begin{table}[htbp]
    \centering
    \caption{Countries most often ranked happiest (WVS, 8 indicators $\times$ 8 samples: each wave and all waves combined)}\label{tab:happiest}
    \input{../tables/main/happiest_countries.tex}
\end{table} % TODO say in the table note (if it's true) that: Looking at all waves combined, Kyrgyzstan–2020 is the happiest country–year according to 3 indicators; Finland, Malaysia, Mexico, Qatar, Vietnam according to other indicators. (give the years). Also add this info in Section 3.6, and say it before the sentence on occurrences

\begin{figure}[htbp]
    \centering
    \caption{Share of \textit{Very happy} people vs.\ GDP per capita (PPP), WVS/EVS 1981--2023}\label{fig:very_happy_gdp}
    \includegraphics[width=.9\textwidth]{../figures/main/wvs_very_happy_vs_gdp_ppp.pdf}
\end{figure}

\begin{figure}[htbp]
    \centering
    \caption{Mean ladder vs.\ GDP per capita (PPP), Gallup, last observation per country}\label{fig:gallup_gdp}
    \includegraphics[width=.9\textwidth]{../figures/main/gallup_ladder_mean_vs_gdp_ppp.pdf}
\end{figure}

\begin{figure}[htbp]
    \centering
    \caption{\textit{Happy} vs.\ GDP per capita (PPP), WVS/EVS 1981--2023}\label{fig:happy_gdp}
    \includegraphics[width=.9\textwidth]{../figures/main/wvs_happy_vs_gdp_ppp.pdf}
\end{figure}

\begin{figure}[htbp]
    \centering
    \caption{\textit{Very Unhappy} vs.\ GDP per capita (PPP), WVS/EVS 1981--2023}\label{fig:very_unhappy_gdp}
    \includegraphics[width=.9\textwidth]{../figures/main/wvs_very_unhappy_vs_gdp_ppp.pdf}
\end{figure}

\begin{figure}[htbp]
    \centering
    \caption{\textit{V.~Happy -- V.~Unhappy} vs.\ GDP per capita (PPP), WVS/EVS 1981--2023}\label{fig:vhmvu_gdp}
    \includegraphics[width=.9\textwidth]{../figures/main/wvs_very_happy_minus_very_unhappy_vs_gdp_ppp.pdf}
\end{figure}

\begin{figure}[htbp]
    \centering
    \caption{\textit{Happiness (mean)} vs.\ GDP per capita (PPP), WVS/EVS 1981--2023}\label{fig:happiness_mean_gdp}
    \includegraphics[width=.9\textwidth]{../figures/main/wvs_happiness_mean_vs_gdp_ppp.pdf}
\end{figure}

\begin{figure}[htbp]
    \centering
    \caption{\textit{Satisfied} vs.\ GDP per capita (PPP), WVS/EVS 1981--2023}\label{fig:satisfied_gdp}
    \includegraphics[width=.9\textwidth]{../figures/main/wvs_satisfied_vs_gdp_ppp.pdf}
\end{figure}

\begin{figure}[htbp]
    \centering
    \caption{\textit{Happy + Satisfied} vs.\ GDP per capita (PPP), WVS/EVS 1981--2023}\label{fig:layard_gdp}
    \includegraphics[width=.9\textwidth]{../figures/main/wvs_happiness_layard_vs_gdp_ppp.pdf}
\end{figure}

\begin{figure}[htbp]
    \centering
    \caption{\textit{Happy} vs.\ GDP per capita (nominal), WVS/EVS wave 7 (2017--2023), weighted by population}\label{fig:happy_gdp_w7_weighted}
    \includegraphics[width=.9\textwidth]{../figures/main/wvs_happy_vs_gdp_wave7_weighted.pdf}
    \begin{minipage}{.9\textwidth}\footnotesize \textit{Note:} Point size and regression weights proportional to population.\end{minipage}
\end{figure}

\begin{figure}[htbp]
    \centering
    \caption{Mean life evaluation vs.\ GDP per capita (PPP), by question variant, Fabre (2025)}\label{fig:fabre_gdp}
    \includegraphics[width=.9\textwidth]{../figures/main/fabre_mean_vs_gdp_ppp.pdf}
    \begin{minipage}{.9\textwidth}\footnotesize \textit{Note:} Mean answer on the native scale of each variant (0--10 or 1--10), ten countries surveyed in 2025. Colors: wording; shapes and line types: scale. Lines: OLS fits by variant.\end{minipage}
\end{figure}


\end{document}
